% Status: bounded-draft-ready. Inhaltlich vollständig, finale Zitation
% nach Quellenreview. Quelle: docs/research/THESIS_KAPITEL_03_METHODIK_DRAFT.md
\section{Daten und Methodik}
\label{ch:methodik}

\subsection{Forschungsdesign: deterministischer Kern vor Interpretation}

Die Arbeit prüft informationelle Effizienz nicht als direkt beobachtbare
Eigenschaft, sondern über drei reproduzierbare Proxy-Tests: Prognosequalität
(H1), Ereignisreaktion (H2) und Wallet-Timing (H3). Diese Operationalisierung
folgt dem Effizienzbegriff der Markteffizienzhypothese \citep{lit_emh_001} und
überträgt ihn auf einen dezentralen Prognosemarkt, dessen Preispfad als
primäre Zeitreihe behandelt wird.

Methodisch gilt eine strikte Trennung: Alle statistischen Kennzahlen werden
deterministisch in Python berechnet, versioniert und als Artefakt (CSV, JSON,
PNG) abgelegt. Sprachmodelle oder Agenten berechnen keine Metriken, sie dürfen
ausschliesslich bereits berechnete, begrenzte Ergebnisse interpretieren, und
auch das erst nach dokumentierter Freigabe und Audit-Protokollierung. Diese
Reihenfolge --- Datenvalidierung, deterministische Analyse, danach Interpretation
--- sichert Nachvollziehbarkeit und schützt vor versteckten Berechnungen im
Fliesstext.

\subsection{Datenbasis}

Die empirische Basis liegt in einer lokalen SQLite-Datenbank
(\texttt{data/thesis.db}). Tabelle~\ref{tab:t1} fasst Quellen, Zeilenzahlen und
bekannte Einschränkungen zusammen.

\begin{table}[ht]
\centering
\caption{Daten- und Quelleninventar (T1)}
\label{tab:t1}
\small
\begin{tabular}{@{}llll@{}}
\toprule
Quelle & Tabelle & Umfang & Rolle / Einschränkung \\
\midrule
Polymarket CLOB/Gamma & \texttt{polymarket\_prices} & 307 Tagespreise & primäre Marktzeitreihe \\
FiveThirtyEight & \texttt{poll\_forecasts} & 245 Zeilen & Vergleichsquelle H1 \\
On-Chain-Wallets & \texttt{whale\_trades} & 25\,113 Zeilen & H3, BUY-only, Filter $\geq$10\,000 USD \\
GDELT & \texttt{sentiment\_scores} & 310 Zeilen & Kontext, kein Kernparameter \\
Ereigniskatalog (Seed) & \texttt{events\_timeline} & 7 kuratierte Ereignisse & H2-Eingabe, vorab fixiert \\
\bottomrule
\end{tabular}
\end{table}

Die Vergleichsquellen sind bewusst restriktiv gewählt. FiveThirtyEight liefert
native Wahrscheinlichkeiten und ist deshalb direkt vergleichbar
\citep{zotero_poly_002}. Für die Kalibrierungsdiagnostik in H1 treten als weitere aufgelöste Vergleichsquellen das unabhängige Umfragemodell Rieke, die 270toWin/JHK-Staatenprognose sowie eine dokumentierte Umfrage-Transformation (Poll-Transform) hinzu, jeweils nur für den paarweisen Brier-Vergleich und nicht als Kernzeitreihe. RealClearPolitics liefert Umfragedurchschnitte, keine
Prognosewahrscheinlichkeiten, diese Quelle bleibt aus allen Wahrscheinlichkeits-
und Kalibrierungsmetriken ausgeschlossen, solange die in der Methodiknote
(\texttt{docs/research/RCP\_TRANSFORMATION.md}) beschriebene Logit-Transformation
nicht freigegeben ist. Die Wallet-Daten tragen die methodisch wichtigste
Datenlimitation: Sie enthalten nur Kauftransaktionen und sind quellseitig ab
10\,000 USD gefiltert, beides bleibt Metadaten-Eigenschaft der Rohquelle und ist
keine analytische Schwellendefinition.

\subsection{Operationalisierung der Hypothesen}

\paragraph{H1 --- Prognosequalität}
H1 bewertet, ob Polymarket im überlappenden Vergleichsfenster bessere
Prognosequalität zeigt als vergleichbare traditionelle
Wahrscheinlichkeitsprognosen. Als Verlustmass dient der Brier-Score
\citep{lit_brier_001}, der paarweise Vergleich der Verlustreihen erfolgt mit dem
Diebold-Mariano-Test \citep{lit_dm_001}, ergänzt um Kalibrierungsdiagnostik,
soweit Stichprobengrösse und Methodik es zulassen. Erlaubt ist ausschliesslich
die Aussage eines Prognosequalitäts-Vergleichs, ein Beweis höherer
Reaktionsgeschwindigkeit, eine allgemeine Marktüberlegenheit oder eine
RCP-Wahrscheinlichkeitsaussage ohne dokumentierte Transformation sind
ausgeschlossen.

\paragraph{H2 --- Ereignisfenster}
H2 prüft, ob sich Polymarket-Wahrscheinlichkeiten um vorab kuratierte
öffentliche Ereignisse in plausibler Richtung und innerhalb fester Tagesfenster
bewegen. Die Ereignisse werden vor der Analyse mit Zeitstempel, Quelle und
erwarteter Richtung fixiert und nach Sichtung der Marktreaktion weder
hinzugefügt noch entfernt. Das Design folgt der Ereignisstudien-Methodik
\citep{lit_eventstudy_001}. Erlaubt ist die Aussage einer täglichen
Ereignisfenster-Reaktion, ein Intraday-Geschwindigkeitsanspruch oder eine
nachträgliche Ereignisauswahl sind ausgeschlossen.

\paragraph{H3 --- Wallet-Timing}
H3 untersucht, ob dataset-relative Wallet-Tiers zeitliche Aktivitätsmuster vor
oder um Polymarket-Preisbewegungen zeigen. Die Tiers werden nicht über eine
willkürliche USD-Schwelle definiert, sondern aus den beobachteten kumulierten
Betragsperzentilen der Wallet-Verteilung abgeleitet. Die Diagnostik kombiniert
deskriptive Lead-Time-Histogramme, Lead-Lag-Korrelationen und Granger-Tests
\citep{lit_granger_001}. Die Granger-Ausgaben werden als prädiktive
Timing-Diagnostik unter Modellannahmen gelesen, nicht als Kausalitätsbeweis,
die Vorsicht gegenüber Insider-Interpretationen stützt sich auf
\citet{zotero_poly_005}. Ausgeschlossen bleiben Aussagen zu Kausalität, privater
Information, Fehlverhalten oder Profitabilität.

\subsection{Reproduzierbarkeit und Guardrails}

Quelldaten liegen portabel in SQLite, analytische Abfragen nutzen zusätzlich
DuckDB. Jede Datenbankschreiboperation wird, wo sinnvoll, validiert
(pydantic/pandera). Der deterministische Kern ist getestet: zum Stand dieses
Entwurfs bestehen 640 Tests. Der Datenzugriff für spätere
Interpretationsschichten ist begrenzt (kein \texttt{SELECT *} ohne
\texttt{LIMIT}, maximal 50 Zeilen pro Abfrage), und jeder spätere
Sprachmodell-Aufruf muss protokolliert werden. So bleibt jede thesis-relevante Kennzahl reproduzierbar und prüfbar.
